\documentclass[11pt]{article}
\usepackage{amsmath, amssymb, amsthm}
\usepackage{geometry}
\usepackage{hyperref}
\usepackage{booktabs}
\usepackage{listings}

\geometry{margin=1.1in}

\newtheorem{lemma}{Lemma}
\newtheorem{fact}{Fact}

\newcommand{\eps}{\varepsilon}
\newcommand{\dt}{\Delta t}
\newcommand{\dx}{\Delta x}
\newcommand{\code}[1]{\texttt{#1}}

% hyperref cannot typeset math or \code in PDF bookmarks; tell it what to use.
\pdfstringdefDisableCommands{%
  \def\theta{theta}\def\Theta{Theta}\def\alpha{alpha}\def\rho{rho}%
  \def\eps{eps}\def\dt{dt}\def\dx{dx}\def\code#1{#1}%
}

\lstset{
  basicstyle=\ttfamily\small,
  columns=fullflexible,
  keepspaces=true,
  breaklines=true,
  frame=single,
}

\title{ETD-$\theta$: Removing the $\eps$-Dependent Quadrature Bias\\ from the ETD
Fokker--Planck Projection}
\author{Optimal Transport / Schr\"odinger Bridge Project\\
Implements \code{projection\_expsemi\_theta.py}; reference scheme is
\code{projection\_expsemi.py}}
\date{}

\begin{document}
\maketitle

\begin{abstract}
The ETD projection in \code{projection\_expsemi.py} integrates the diffusion term
\emph{exactly}, yet its accuracy still degrades as $\eps$ grows: measured on the FP
solution error, its order falls from $2.00$ at $\eps=10^{-2}$ to $1.28$ at $\eps=10$,
losing a factor $60$ in absolute accuracy. This note shows the cause is not the
semigroup treatment but a \emph{source-term quadrature bias}: ETD1 samples $D_xm$ at
the time midpoint, while the exponential weight it is integrated against has its
centroid at $\theta(\alpha)>1/2$. Correcting the sample point --- by interpolating
between the two $m$ rows the staggered grid already provides, with weight
$\omega=\theta(\alpha)-1/2$ --- restores second order uniformly in $\eps$ with a
constant flat across four decades, at the cost of one extra precomputed array. The
correction preserves the tridiagonal-in-$t$ structure, so the existing batched Thomas
sweep still applies, and it degenerates to ETD1 identically at $\omega=0$. Sections
\ref{sec:lastrow} and \ref{sec:open} record two things that did \emph{not} go as
expected: the final time row forces a specific extrapolation choice (a linear one
destroys the band structure), and the end-to-end LADMM confirmation at large $\eps$ is
currently inconclusive for reasons unrelated to the scheme.
\end{abstract}

\tableofcontents

\section{Setup and notation}

Throughout, the FP constraint is
\begin{equation}
  \partial_t\rho + \partial_x m = \eps\,\partial_{xx}\rho,
  \label{eq:fp}
\end{equation}
discretised on the staggered grid of \code{operators.py}: $\rho$ at $(n_t-1)\times n_x$
(time-interior nodes, space cell-centres), $m$ at $n_t\times(n_x-1)$ (time cell-centres,
space-interior nodes), and the residual on the $\phi$ grid $n_t\times n_x$.

Every scheme here is diagonal in the DCT-$x$ basis. Writing $\lambda_j$ for the discrete
Neumann Laplacian eigenvalues of \code{grid.py},
\begin{equation}
  \lambda_j = \frac{2-2\cos(j\pi/n_x)}{\dx^2},
  \qquad
  \alpha_j := \eps\,\lambda_j\,\dt ,
  \label{eq:alpha}
\end{equation}
and hatting DCT-$x$ coefficients, ETD1 reads, per mode $j$ and time row
$k=0,\dots,n_t-1$,
\begin{equation}
  \frac{\hat\rho_{k+1} - c_j\hat\rho_k}{\dt} + \varphi_j\,\hat g_k = 0,
  \qquad
  c_j = e^{-\alpha_j},\quad
  \varphi_j = \frac{1-c_j}{\alpha_j},
  \label{eq:etd1}
\end{equation}
where $g := D_xm$ and $D_x$ is \code{deriv\_x\_at\_phi}. The boundary rows use
$\hat\rho_0=\hat\rho_0^{\mathrm{BC}}$ and $\hat\rho_{n_t}=\hat\rho_1^{\mathrm{BC}}$.
Note $\hat g_k$ is $m$ at the time \emph{midpoint} $t_{k+1/2}$, since that is where the
staggered grid places $m$.

\section{The error is a quadrature bias, not a semigroup error}

The factor $c_j=e^{-\alpha_j}$ in \eqref{eq:etd1} is the \emph{exact} semigroup factor
for the space-discrete problem, so the homogeneous part of ETD1 carries no error at any
$\eps$. This is worth stating sharply because it is verifiable: feeding ETD1 the
semi-discrete pure-diffusion solution
$\rho(t,x_i)=1+a\cos(n\pi x_i)e^{-\eps\lambda_nt}$ with $m\equiv0$ produces a residual of
$10^{-14}$ for every $\eps$ from $10^{-3}$ to $10$, whereas Crank--Nicolson on the same
data produces $1.6\times10^{-9}$ up to $1.2\times10^{1}$ over that range.

The error lives entirely in the source term. Integrating \eqref{eq:fp} exactly over one
step gives
\begin{equation}
  \hat\rho_{k+1}
  = c_j\hat\rho_k
  - \int_{t_k}^{t_k+\dt} e^{-\eps\lambda_j(t_{k+1}-s)}\,\hat g(s)\,ds .
  \label{eq:exactstep}
\end{equation}
ETD1 approximates the integral by $\dt\,\varphi_j\,\hat g(t_{k+1/2})$. The weight
\emph{integral} is exact, since
$\int_{t_k}^{t_{k+1}}e^{-\eps\lambda_j(t_{k+1}-s)}ds=\dt\varphi_j$. What is wrong is the
sample point: the weight
\begin{equation}
  w(s) = e^{-\eps\lambda_j(t_{k+1}-s)}
\end{equation}
is not symmetric about the midpoint --- it increases monotonically towards
$s=t_{k+1}$. Expanding $\hat g$ about the midpoint $t_m := t_{k+1/2}$,
\begin{equation}
  \int w(s)\bigl[\hat g(s)-\hat g(t_m)\bigr]ds
  = \hat g'(t_m)\int w(s)(s-t_m)\,ds + O(\dt^3),
\end{equation}
and the first-order term survives precisely because $t_m$ is not the weight's centroid.

\begin{fact}[the bias]
With $\theta(\alpha)$ the $w$-weighted centroid of $[t_k,t_{k+1}]$ expressed as a
fraction of $\dt$, the leading ETD1 error per step is
\begin{equation}
  \dt\,\varphi_j\cdot\dt\,\bigl(\theta(\alpha_j)-\tfrac12\bigr)\,\hat g'(t_m).
  \label{eq:bias}
\end{equation}
It vanishes at $\eps=0$ (flat weight, centroid $=$ midpoint) and saturates as
$\alpha\to\infty$. This is the $\eps$-dependent error constant.
\end{fact}

\section{The centroid weight $\theta(\alpha)$}
\label{sec:theta}

Substitute $u=(s-t_k)/\dt$, so $u\in[0,1]$ and $w(u)=e^{-\alpha(1-u)}$. The two moments
are
\begin{align}
  \int_0^1 e^{-\alpha(1-u)}\,du       &= \frac{1-e^{-\alpha}}{\alpha}
                                       =: \varphi(\alpha), \label{eq:phi}\\
  \int_0^1 u\,e^{-\alpha(1-u)}\,du    &= \frac{(\alpha-1)+e^{-\alpha}}{\alpha^2}
                                       =: \psi(\alpha), \label{eq:psi}
\end{align}
the second by parts. Their ratio is the centroid:
\begin{equation}
  \boxed{\;
  \theta(\alpha) = \frac{\psi(\alpha)}{\varphi(\alpha)}
                 = \frac{(\alpha-1)+e^{-\alpha}}{\alpha\bigl(1-e^{-\alpha}\bigr)}\;}
  \label{eq:theta}
\end{equation}

\paragraph{Limits.} Expanding \eqref{eq:phi}--\eqref{eq:psi},
$\psi = \tfrac12-\tfrac{\alpha}{6}+\tfrac{\alpha^2}{24}-\tfrac{\alpha^3}{120}+\cdots$
and $1/\varphi = 1+\tfrac{\alpha}{2}+\tfrac{\alpha^2}{12}+0\cdot\alpha^3+\cdots$, so
\begin{equation}
  \theta(\alpha) = \frac12 + \frac{\alpha}{12} - \frac{\alpha^3}{720} + O(\alpha^5),
  \label{eq:thetaseries}
\end{equation}
with the $\alpha^2$ and $\alpha^4$ terms vanishing. As $\alpha\to\infty$,
$e^{-\alpha}\to0$ and $\theta\to(\alpha-1)/\alpha = 1-1/\alpha\to1$: the weight collapses
onto the right endpoint. Hence $\theta\in[\tfrac12,1)$, monotone increasing, and the
midpoint is correct only to $O(\alpha)$.

\paragraph{The interpolation weight.} Define
\begin{equation}
  \omega(\alpha) := \theta(\alpha)-\tfrac12 \;\in\; [0,\tfrac12).
  \label{eq:omega}
\end{equation}
The target sample point $t_k+\theta\dt$ lies in $[t_{k+1/2},t_{k+1}]$. The two nearest
$m$ nodes are $t_{k+1/2}$ and $t_{k+3/2}$, and the target sits at fractional position
$\theta+\tfrac12\in[1,1.5]$ between them --- always \emph{inside} that pair, never an
extrapolation. Linear interpolation therefore gives
\begin{equation}
  \hat g(t_k+\theta\dt) \;\approx\; (1-\omega)\,\hat g_k + \omega\,\hat g_{k+1}.
  \label{eq:interp}
\end{equation}
The first-order bias \eqref{eq:bias} is removed by construction; what remains is the
$O(\dt^2)$ interpolation error times $\int w=\dt\varphi$, and since $\varphi\le1$ for
all $\alpha$, the remaining constant is bounded uniformly in $\eps$.

\begin{lemma}[relation to ETD2]
\eqref{eq:interp} is ETD2's quadrature. ETD2 integrates the linear interpolant of $\hat
g$ through two \emph{integer}-node values exactly against $w$, giving
$\varphi\bigl[(1-\theta)\hat g(t_k)+\theta\hat g(t_{k+1})\bigr]$ with $\theta=\psi/\varphi$
as in \eqref{eq:theta} --- the same rule, evaluated on a different grid.
\end{lemma}

The practical consequence is that the $n_t+1$ integer-node relayout assumed by the
(unimplemented) \code{cfg\_ladmm\_gaussian\_expsemi\_etd2.m} is unnecessary: the
half-integer nodes already carry enough information to achieve ETD2's order.

\section{The scheme and its normal equations}

ETD-$\theta$ replaces \eqref{eq:etd1} by
\begin{equation}
  \bigl[R(\rho,m)\bigr]_k
  := \frac{\hat\rho_{k+1}-c_j\hat\rho_k}{\dt}
   + \varphi_j\bigl[\Theta\hat g\bigr]_k = 0,
  \label{eq:etdtheta}
\end{equation}
with $\Theta$ the time-only operator implementing \eqref{eq:interp}. Only the $m$ block
changes; the $\rho$ block is bit-for-bit ETD1's.

\subsection{$\Theta$ and its adjoint}

As implemented (\code{\_apply\_theta}, \code{\_apply\_theta\_adj}), with $\omega$
per-mode:
\begin{equation}
  [\Theta v]_k =
  \begin{cases}
    (1-\omega)v_k + \omega v_{k+1}, & 0\le k\le n_t-2,\\[2pt]
    v_{n_t-1},                      & k=n_t-1,
  \end{cases}
  \label{eq:Theta}
\end{equation}
so as a matrix $\Theta$ is upper bidiagonal with diagonal
$(1-\omega,\dots,1-\omega,1)$. Accumulating $[\Theta^*v]_l=\sum_k\Theta_{kl}v_k$ column
by column,
\begin{equation}
  [\Theta^*v]_l =
  \begin{cases}
    (1-\omega)v_0,                        & l=0,\\[2pt]
    (1-\omega)v_l + \omega v_{l-1},       & 1\le l\le n_t-2,\\[2pt]
    v_{n_t-1} + \omega v_{n_t-2},         & l=n_t-1.
  \end{cases}
  \label{eq:ThetaStar}
\end{equation}

\subsection{The per-mode Gram matrix}

Since $\Theta$ acts on time and $D_x$ on space, they commute, and $R_m=\varphi\Theta D_x$
gives $R_mR_m^*=\varphi^2\Theta(D_xD_x^*)\Theta^*=\varphi^2\lambda\,\Theta\Theta^*$. With
$R_\rho R_\rho^*$ unchanged from ETD1, the per-mode system is
\begin{equation}
  T_j = R_\rho R_\rho^* + \varphi_j^2\lambda_j\,\Theta\Theta^* .
  \label{eq:gram}
\end{equation}
Both terms are tridiagonal, hence so is $T_j$. Writing $\Gamma:=\varphi_j^2\lambda_j$,
the entries implemented in \code{precomp\_expsemi\_theta\_proj} are

\begin{center}
\begin{tabular}{@{}lll@{}}
\toprule
row & main diagonal & off-diagonal (row $k$, $k{+}1$) \\
\midrule
$k=0$              & $1/\dt^2 + \Gamma[(1-\omega)^2+\omega^2]$
                   & $-c/\dt^2 + \Gamma\,\omega(1-\omega)$ \\
$1\le k\le n_t-3$  & $(1+c^2)/\dt^2 + \Gamma[(1-\omega)^2+\omega^2]$
                   & $-c/\dt^2 + \Gamma\,\omega(1-\omega)$ \\
$k=n_t-2$          & $(1+c^2)/\dt^2 + \Gamma[(1-\omega)^2+\omega^2]$
                   & $-c/\dt^2 + \Gamma\,\omega$ \\
$k=n_t-1$          & $c^2/\dt^2 + \Gamma$
                   & --- \\
\bottomrule
\end{tabular}
\end{center}

\noindent
The two irregular entries both come from $\Theta$'s final row being a lone $1$: its Gram
diagonal is $1$ rather than $(1-\omega)^2+\omega^2$, and
$(\Theta\Theta^*)_{n_t-2,n_t-1}=\omega$ rather than $\omega(1-\omega)$. Setting
$\omega=0$ collapses the table to $\mathrm{diag}=(1+c^2)/\dt^2+\varphi^2\lambda$ with
constant off-diagonal $-c/\dt^2$, which is exactly \code{precomp\_expsemi\_proj}: the
degeneracy to ETD1 is structural, not approximate.

Because the off-diagonal is row-varying in one place, the module uses
\code{thomas\_batch\_precomp\_general} / \code{thomas\_batch\_solve\_general} from
\code{projection\_nu.py} rather than \code{projection.py}'s constant-off-diagonal pair.
These reduce to the constant version when \code{sub} is uniform, so nothing is lost.

\subsection{The correction step}

The projection is $x_{\mathrm{out}}=x_{\mathrm{in}}-R^*T^{-1}Rx_{\mathrm{in}}$. Since
$R_\rho$ is untouched, the $\rho$ correction is ETD1's verbatim. For $m$,
$R_m^*=D_x^*\Theta^*\varphi$, so $\Theta^*$ is applied in DCT space (where $\omega$ is
diagonal) before the inverse transform:
\begin{lstlisting}
p_weighted = idct(_apply_theta_adj(phi_vals * p_hat, omega), axis=1)
mx_out     = psi + ops.deriv_x_at_m(p_weighted)
\end{lstlisting}
The absent minus sign is correct: \code{deriv\_x\_at\_m} is the \emph{negated} adjoint of
\code{deriv\_x\_at\_phi}, as a two-point difference must be, and ETD1 relies on the same
cancellation.

\paragraph{The DC mode.} At $j=0$, $\lambda_0=0$ so $\alpha_0=0$, giving $c=1$,
$\varphi=1$, $\omega=0$, and the entire $m$ block drops out of \eqref{eq:gram}. $T_0$ is
therefore ETD1's singular Neumann-in-$t$ Laplacian, handled by the same DCT-in-$t$ gauge
fix, unchanged.

\section{The final time row forces the extrapolation choice}
\label{sec:lastrow}

Row $k=n_t-1$ of \eqref{eq:interp} needs $\hat g_{n_t}$, which does not exist: $m$ has
$n_t$ rows, the last at $t=1-\dt/2$. The obvious choice is linear extrapolation,
$\hat g_{n_t}:=2\hat g_{n_t-1}-\hat g_{n_t-2}$, which keeps that row second-order
consistent. \textbf{It must not be used.}

With linear extrapolation, $\Theta$'s final row spans columns $\{n_t-2,n_t-1\}$. That row
then overlaps row $n_t-3$, which spans $\{n_t-3,n_t-2\}$, producing
\begin{equation}
  (\Theta\Theta^*)_{n_t-3,\,n_t-1}
  = \omega\cdot(-\omega) = -\omega^2 \neq 0 ,
\end{equation}
a distance-2 entry. $\Theta\Theta^*$ becomes bandwidth 2, \eqref{eq:gram} is no longer
tridiagonal, and the Thomas sweep silently solves the wrong system. This was not caught
by hand --- it was caught by check \#5 of \code{scripts/validate\_expsemi\_theta.py},
which reported post-projection residual ratios of $3\times10^{-2}$ instead of
$10^{-15}$, together with a $5\times10^{-3}$ idempotence drift.

The argument generalises: \emph{any} two-point final row touching column $n_t-2$ overlaps
row $n_t-3$. A tridiagonal $RR^*$ therefore requires the final row to touch column
$n_t-1$ alone, i.e. constant extrapolation $\hat g_{n_t}:=\hat g_{n_t-1}$, which collapses
that row to the plain midpoint sample --- \eqref{eq:Theta}.

Two consolations. First, $\Theta$ is then upper triangular with
$\det\Theta=(1-\omega)^{n_t-1}\neq0$ for $\omega<1$, so $\Theta\Theta^*$ is SPD by
inspection, with no continuity argument needed. Second, the accuracy cost is nil:
Table~\ref{tab:lastrow} compares the two rules and they agree to every printed digit. The
reason is structural --- the final residual row determines $\hat\rho_{n_t}$, which is a
prescribed boundary condition rather than an unknown.

\begin{table}[htbp]
\centering
\caption{Last-row rule, relative $L^2$ FP solution error, rough $\sigma=0.05$ data,
$n_t=n_x=n$. ``const'' is what ships; ``linear'' is the bandwidth-2 variant.}
\label{tab:lastrow}
\begin{tabular}{@{}llrrrrr@{}}
\toprule
$\eps$ & rule & $n{=}32$ & $n{=}64$ & $n{=}128$ & $n{=}256$ & $n{=}512$ \\
\midrule
$0.1$ & midpoint (ETD1) & 1.818e-2 & 4.755e-3 & 1.212e-3 & 3.046e-4 & 7.624e-5 \\
      & const           & 1.098e-2 & 2.780e-3 & 7.074e-4 & 1.789e-4 & 4.499e-5 \\
      & linear          & 1.098e-2 & 2.780e-3 & 7.074e-4 & 1.789e-4 & 4.499e-5 \\
\midrule
$1$   & midpoint (ETD1) & 4.399e-2 & 1.661e-2 & 5.789e-3 & 1.811e-3 & 5.100e-4 \\
      & const           & 9.580e-3 & 2.403e-3 & 6.262e-4 & 1.652e-4 & 4.327e-5 \\
      & linear          & 9.580e-3 & 2.403e-3 & 6.262e-4 & 1.652e-4 & 4.327e-5 \\
\midrule
$10$  & midpoint (ETD1) & 6.512e-2 & 2.967e-2 & 1.320e-2 & 5.614e-3 & 2.234e-3 \\
      & const           & 9.531e-3 & 2.318e-3 & 5.752e-4 & 1.444e-4 & 3.687e-5 \\
      & linear          & 9.531e-3 & 2.318e-3 & 5.752e-4 & 1.444e-4 & 3.687e-5 \\
\bottomrule
\end{tabular}
\end{table}

\section{Evaluating $\theta$ in floating point}

The direct form \eqref{eq:theta} is catastrophically cancelling at small $\alpha$: its
numerator $(\alpha-1)+e^{-\alpha}\sim\alpha^2/2$ is a difference of $O(1)$ quantities.
The damage is invisible in $\theta$ and severe in $\omega$, because $\theta\approx1/2$
dominates its own relative error while $\omega=\theta-1/2$ does not. At $\alpha=10^{-4}$
the direct form is already wrong in the fourth significant digit of $\omega$ (relative
error $1.5\times10^{-4}$) while $\theta$ still looks accurate to $10^{-12}$.

Both branches were compared against adaptive quadrature of
\eqref{eq:phi}--\eqref{eq:psi} over $\alpha\in[10^{-6},3]$. They cross near
$\alpha=0.02$, where each holds about $10^{-10}$ relative error on $\omega$; the module
switches there, using the two-term series \eqref{eq:thetaseries} below and the direct
form above. For large $\alpha$, $e^{-\alpha}$ underflows harmlessly and \eqref{eq:theta}
is well conditioned.

The lesson generalises to the validation: \textbf{check the error on $\omega$, not on
$\theta$}. The first version of check \#1 measured $\theta$ and passed a cutoff that was
two orders of magnitude too low.

\section{Files}

\begin{center}
\begin{tabular}{@{}lll@{}}
\toprule
file & status & what \\
\midrule
\code{optimal\_transport/projection\_expsemi\_theta.py} & new
  & the scheme: $\theta,\omega$, $\Theta,\Theta^*$, precomp, projection \\
\code{scripts/validate\_expsemi\_theta.py} & new
  & the six checks of \S\ref{sec:validation} \\
\code{scripts/compare\_etd1\_vs\_theta.py} & new
  & end-to-end LADMM comparison (\S\ref{sec:open}) \\
\code{scripts/study\_utils.py} & edited
  & \code{run\_study} gained \code{precomp}/\code{projection} arguments \\
\code{optimal\_transport/projection\_expsemi.py} & \textbf{untouched}
  & remains the ETD1 reference \\
\code{optimal\_transport/projection.py} & \textbf{untouched}
  & remains the CN reference \\
\bottomrule
\end{tabular}
\end{center}

The \code{run\_study} change defaults to the ETD1 pair, so every existing caller
(\code{grid\_convergence.py}, \code{multi\_eps\_study.py}, \dots) is unaffected. The new
scheme plugs into the \code{projection}/\code{projection\_state} seam that
\code{discretize\_then\_optimize} already exposed, so \code{pipeline.py} needed no change
at all.

\section{Validation}
\label{sec:validation}

\code{scripts/validate\_expsemi\_theta.py} runs six checks over
$\eps\in\{0,10^{-8},10^{-3},0.1,1,10,10^3\}$ at $n_t=32$, $n_x=24$. Nothing trusts the
hand algebra of \S\ref{sec:theta}--\S\ref{sec:lastrow}: each piece is checked against an
independently constructed reference.

\begin{enumerate}
  \item $\theta,\omega$ in range and monotone; $\omega$ matches adaptive quadrature of
        \eqref{eq:phi}--\eqref{eq:psi} to $<10^{-9}$ relative; the $\alpha\to\infty$
        asymptote $\theta\sim1-1/\alpha$ holds to $O(\alpha^{-2})$.
  \item $\langle\Theta g,v\rangle=\langle g,\Theta^*v\rangle$ to $2.7\times10^{-15}$, and
        \code{\_apply\_theta} agrees with a densely built $\Theta$ to
        $4.4\times10^{-16}$, for $\omega\in\{0,10^{-6},0.05,0.2,0.4999\}$.
  \item The closed-form $(\mathrm{main},\mathrm{sub})$ of \eqref{eq:gram} matches a dense
        $R_\rho R_\rho^*+\varphi^2\lambda\,\Theta\Theta^*$ built per mode to
        $3.4\times10^{-18}$ relative; every per-mode $T_j$ is verified tridiagonal and
        SPD (worst $\lambda_{\min}/\lambda_{\max}=4.4\times10^{-10}$ over all modes and
        all $\eps$).
  \item At $\eps=0$, ETD-$\theta$ reproduces \code{proj\_fokker\_planck\_expsemi}
        \emph{exactly} (max relative difference $0$).
  \item \textbf{Feasibility.} The ETD-$\theta$ residual of the projected point, recomputed
        with a dense $\Theta$ rather than the module's own helper, drops from
        $O(10^3)$ to $O(10^{-12})$ --- worst after/before ratio $8.2\times10^{-15}$
        across the whole $\eps$ sweep. This is the decisive check: a projection built on
        a $T$ that is not the true Gram $RR^*$ cannot land on the constraint set, so
        \#5 would fail if \emph{any} of \eqref{eq:Theta}, \eqref{eq:ThetaStar} or
        \eqref{eq:gram} were wrong. It is what caught the bandwidth-2 error of
        \S\ref{sec:lastrow}.
  \item Idempotence: $P(P(x))=P(x)$ to $0$ (exactly).
\end{enumerate}

\section{Measured effect}

The scheme-level measurement marches each scheme forward in time from $\rho(0)$ using the
exact source $g=D_xm$ and compares against the exact solution --- no ADMM, no projection,
so what is measured is purely FP discretisation error. Data is a manufactured solution
with a Gaussian-decaying spectrum at the project's own $\sigma=0.05$, which has genuine
high-frequency content (31 $x$-modes in 1D, $11\times11$ in 2D).

Solution error, not residual norm, is the right measure here: for a stiff mode the
residual carries a factor $1/\dt$ that the inverse operator removes, so a residual-fitted
order understates every scheme and does so unevenly.

\begin{table}[htbp]
\centering
\caption{Relative $L^2$ FP solution error at the finest grid, with fitted order $p$.
1D: $n_t=n_x=512$. 2D: $n_t=n_x=n_y=96$, built with the project's own
\code{operators\_2d}/\code{grid\_2d}.}
\label{tab:main}
\begin{tabular}{@{}lrrrrrr@{}}
\toprule
& \multicolumn{2}{c}{CN} & \multicolumn{2}{c}{ETD1}
& \multicolumn{2}{c}{ETD-$\theta$} \\
\cmidrule(lr){2-3}\cmidrule(lr){4-5}\cmidrule(lr){6-7}
$\eps$ & error & $p$ & error & $p$ & error & $p$ \\
\midrule
\multicolumn{7}{@{}l}{\emph{1D}}\\
$0.01$ & 4.402e-5 & 2.00 & 4.609e-5 & 2.00 & \textbf{4.163e-5} & 2.00 \\
$0.1$  & 5.093e-5 & 2.00 & 7.624e-5 & 2.00 & \textbf{4.499e-5} & 1.99 \\
$1$    & 5.279e-5 & 2.00 & 5.100e-4 & 1.75 & \textbf{4.327e-5} & 1.93 \\
$10$   & 5.312e-5 & 2.02 & 2.234e-3 & 1.28 & \textbf{3.687e-5} & 1.98 \\
\midrule
\multicolumn{7}{@{}l}{\emph{2D}}\\
$0.01$ & 1.641e-3 & 2.01 & 1.730e-3 & 2.01 & \textbf{1.512e-3} & 2.01 \\
$0.1$  & 2.011e-3 & 2.03 & 3.310e-3 & 1.93 & \textbf{1.695e-3} & 1.96 \\
$1$    & 2.066e-3 & 2.06 & 1.491e-2 & 1.35 & \textbf{1.428e-3} & 1.98 \\
$10$   & 2.167e-3 & 2.19 & 2.563e-2 & 1.08 & \textbf{1.358e-3} & 2.02 \\
\bottomrule
\end{tabular}
\end{table}

Three readings of Table~\ref{tab:main}.

\begin{enumerate}
  \item \textbf{ETD-$\theta$ holds $p\approx2$ at every $\eps$ with a constant flat in
        $\eps$} (1D: 4.2, 4.5, 4.3, $3.7\times10^{-5}$ across four decades). It is the
        best scheme at every $\eps$ tested and never worse anywhere. Against ETD1 it
        gains a factor $60$ at $\eps=10$ and $12$ at $\eps=1$ in 1D, $19$ and $10$ in 2D.
  \item \textbf{ETD1 degrades exactly as \eqref{eq:bias} predicts}, $p$ falling
        $2.00\to1.75\to1.28$.
  \item \textbf{CN's accuracy is not the problem.} Its order and constant are both stable
        in $\eps$. CN's defect is elsewhere --- $R_{\mathrm{CN}}(\infty)=-1$, so at
        $\eps=10$ with $\sigma=0.05$ data $99.2\%$ of modes have $R<0$ and one step
        returns $\min\rho=-3.61$ where the exact solution has $\min\rho=+0.50$. That is
        a positivity failure, not an accuracy failure, and ETD-$\theta$ inherits ETD1's
        exact semigroup so it does not share it.
\end{enumerate}

\paragraph{Higher dimensions.} The correction is purely temporal, and time is
one-dimensional in every spatial dimension: with
$\alpha=\eps(\lambda_x+\lambda_y)\dt$ still a scalar per mode, $\theta,\omega,\Theta$ are
the identical scalar functions evaluated on an $(n_x,n_y)$ array, and the two-row $m$
stencil keeps $RR^*$ tridiagonal in $t$. It also matters \emph{more} as dimension grows:
the fraction of modes with $\alpha>2$ at $n=64$, $\eps=0.01$ is $30\%$ in 1D, $66\%$ in
2D, $87\%$ in 3D, since $\alpha_{\max}=4d\eps n$ is linear in $d$ and $\lambda$ is a sum.

\section{Open items}
\label{sec:open}

\paragraph{On the actual Gaussian SB problem the gain is a few percent, not
$12$--$60\times$.} This is the central caveat and it qualifies
Table~\ref{tab:main}. Plugging the \emph{reference} SB solution into each
scheme's discrete FP residual --- which, since the reference solves the
continuous equation, measures that scheme's truncation error on the real
problem --- gives:

\begin{center}
\begin{tabular}{@{}lrrrr@{}}
\toprule
$\eps$ & $n$ & $\|R_{\mathrm{ETD1}}\|$ & $\|R_{\mathrm{ETD}\text{-}\theta}\|$ & gain \\
\midrule
$0.001$ & 128 & 1.898e-2 & 1.887e-2 & $1.01\times$ \\
$0.01$  & 128 & 3.622e-2 & 3.550e-2 & $1.02\times$ \\
$0.1$   & 128 & 6.481e-1 & 6.129e-1 & $1.06\times$ \\
$1$     & 128 & 1.001e1  & 1.002e1  & $1.00\times$ \\
\bottomrule
\end{tabular}
\end{center}

The reason is structural, and it is the honest limit of this work.
ETD1 loses \emph{order} only when the solution carries significant content in
modes with $\alpha=\eps\lambda\dt\gtrsim1$; for modes it resolves, ETD1 and
ETD-$\theta$ are both second order and differ by a bounded constant. The
manufactured solutions of Table~\ref{tab:main} were built with a
$\sigma=0.05$ spectrum across 31 $x$-modes \emph{and} a momentum field
$m\propto\eps k^2g$, so their $D_xm$ is large and fast-varying exactly in the
stiff modes the correction targets. The Gaussian SB solution is not like
that: it is smooth, and at large $\eps$ it becomes \emph{smoother} (more
nearly uniform), so it has almost no content where $\alpha\gtrsim1$.

\textbf{So ETD-$\theta$ removes a real, correctly identified $\eps$-dependent
error term, but that term is not the dominant one for this problem.} The
original question --- why the error grows with $\eps$ here --- remains open.
Note in particular that at $\eps=0.1$ the fitted order degrades to $\approx1$
for \emph{both} schemes equally, which rules the quadrature bias out as the
cause and points at a separate $\eps$-dependent term not yet identified
(candidates: the spatial treatment of $\eps\partial_{xx}$, or the reference's
own boundary mismatch).

\paragraph{The reference is degenerate at large $\eps$.} At $\eps=1$ the
residual above does not decrease under refinement at all ($8.79\to9.97\to
10.01$ for $n=32,64,128$): the reference does not satisfy the discrete FP
equation, so nothing measured against it at that $\eps$ means anything.
\code{sinkhorn\_fit} applies the heat semigroup over the \emph{whole}
interval $T=1$, and $e^{-\eps\pi^2}=5.2\times10^{-5}$ at $\eps=1$ annihilates
every non-constant mode, so the iteration reaches a trivial fixed point in 2
iterations. This is adjacent to the known
\code{SINKHORN\_VAREPS\_THRESHOLD} issue and needs fixing before any
large-$\eps$ claim, for any scheme, can be made.

\paragraph{The end-to-end LADMM comparison is correspondingly uninformative.}
\code{scripts/compare\_etd1\_vs\_theta.py} shows ETD-$\theta$ ahead by
$1.01$--$1.08\times$, consistent with the residual table above, with absolute
errors of $O(1)$--$O(10)$ at $\eps\ge1$ against the degenerate reference. At
$\eps=0.1$, where the reference is most trustworthy, the gain does grow
monotonically with resolution ($1.03\to1.05\to1.08$) --- the right direction,
but a weak signal.

\paragraph{Where ETD-$\theta$ should pay off.} A problem retaining genuine
fine structure at large $\eps$ --- bimodal or rough marginals, cf.\
\code{matlab/discretize\_first/2d/experiments/test\_expsemi\_bimodal.m} ---
puts solution content where $\alpha\gtrsim1$ and is the natural test. This has
not been run.

\paragraph{Also not done.}
\begin{itemize}
  \item The \code{\_nu} (non-uniform time grid) variant. It matters more there, since
        grading shrinks $\dt$ and raises the local $\alpha$. \code{projection\_expsemi\_nu.py}
        builds its system by numerical probing rather than hand algebra, so the
        correction should be absorbed with no new derivation.
  \item The 2D/3D port. Python 2D currently has only the banded/CN projection, so this
        means porting ETD anyway; it should be ported directly as ETD-$\theta$, with
        \code{matlab/shared/2d/projection/proj\_fokker\_planck\_expsemi.m} as the
        structural reference. One caution: do not materialise the lookahead
        $\hat g_{k+1}$ as a full $(n_t,n_x,n_y)$ array --- fuse it into the existing
        $t$-sweep.
  \item An adjoint check of the full $R$/$R^*$ pair against the LADMM's own norms. Checks
        \#2 and \#5 cover $\Theta$ and the Gram, but the assembled operator is only
        checked implicitly through feasibility.
\end{itemize}

\paragraph{Worth an independent look.} Equations \eqref{eq:psi} (integration by parts),
\eqref{eq:thetaseries} (the vanishing $\alpha^2$ term, which is what makes the two-term
series worth using), \eqref{eq:ThetaStar} (the adjoint's boundary rows), and the
$k=n_t-2$ off-diagonal $\Gamma\omega$ in the table of \S4.2.

\end{document}
